\chapter{The Universal Regime Map}
\label{ch:regime_map}

\begin{objectivebox}
\begin{itemize}
    \item Classify the four universal operating regimes (Linear, Nonlinear, Yield, Ruptured) of the spatial LC network.
    \item Define the universal dimensionless control parameter $r = A/A_c$ for cross-domain analysis.
    \item Understand how identical saturation mathematics map exactly to semiconductor avalanche breakdowns, black hole event horizons, and Schwinger pair-production domains.
\end{itemize}
\end{objectivebox}

Every physical domain in the AVE framework reduces to a single dimensionless control parameter $r = A/A_c$, where $A$ is the local strain amplitude and $A_c$ is the domain-specific critical threshold derived from the four axioms. The saturation operator $S(r) = \sqrt{1 - r^2}$ (Axiom~4) changes character at well-defined boundaries, partitioning all of physics into four universal regimes.

This chapter establishes the regime classification as a \textbf{prerequisite gate}: no domain analysis proceeds without first identifying its regime and verifying the appropriate equation set.

\section{The Four Universal Regimes}

\begin{resultbox}{Universal Regime Classification}
\begin{equation}
    S(r) = \sqrt{1 - r^2}, \qquad r \equiv \frac{A}{A_c}
\end{equation}
\end{resultbox}

\begin{center}
\small
\renewcommand{\arraystretch}{1.2}
\begin{tabularx}{\textwidth}{@{}c l c c l X@{}}
\toprule
\textbf{Regime} & \textbf{Name} & \textbf{$r$ range} & \textbf{$S$ range} & \textbf{Physics} & \textbf{EE Analog} \\ \midrule
I   & Linear     & $r < \sqrt{2\alpha}$        & $S > 0.993$       & Standard equations  & Small-Signal \\
II  & Nonlinear  & $\sqrt{2\alpha} \leq r < \frac{\sqrt{3}}{2}$ & $0.500 < S < 0.993$ & Axiom~4 active  & Large-Signal \\
III & Yield      & $\frac{\sqrt{3}}{2} \leq r < 1$ & $0 < S < 0.500$   & Phase transition  & Avalanche ($M \geq 2$) \\
IV  & Ruptured   & $r \geq 1.0$       & $S = 0$           & Topology destroyed  & Breakdown ($M \to \infty$) \\
\bottomrule
\end{tabularx}
\end{center}

\noindent The boundary values are \textbf{derived from first principles}, not chosen by convention:

\begin{itemize}
    \item \textbf{$r_1 = \sqrt{2\alpha} \approx 0.1208$ (Small-Signal Limit):} The perturbative expansion (Eq.~\ref{eq:S_taylor}) gives a first-order correction $\Delta S = r^2/2$. The lattice's own self-coupling is $\alpha$ (the fine-structure constant = packing fraction / $8\pi$). When $\Delta S = \alpha$:
    \begin{equation}
        \frac{r^2}{2} = \alpha \qquad \Longrightarrow \qquad r_1 = \sqrt{2\alpha} \approx 0.1208
    \end{equation}
    Below $r_1$, Axiom~4 corrections are \textit{sub-$\alpha$}: smaller than the lattice's own coupling strength and therefore physically unresolvable. This is the exact analog of small-signal linearization in semiconductor analysis, where perturbations smaller than $V_T = kT/q$ are absorbed by thermal noise.

    \item \textbf{$r_2 = \sqrt{3}/2 \approx 0.8660$ (Avalanche Onset):} The quality factor $Q(r) = 1/S(r)$ measures energy trapping efficiency. A mode with $Q \geq \ell$ stores more energy per cycle than it radiates. The minimum non-trivial multipole is $\ell = 2$ (dipole cannot radiate gravitational waves; monopole cannot radiate at all). The yield regime begins when $Q$ first reaches this minimum:
    \begin{equation}
        Q(r) = \frac{1}{S(r)} = \ell_{\min} = 2 \qquad \Longrightarrow \qquad S = \frac{1}{2}, \quad r_2 = \sqrt{1 - \frac{1}{4}} = \frac{\sqrt{3}}{2}
    \end{equation}
    In semiconductor terms, this is the onset of \textbf{avalanche multiplication}: the Miller factor $M = 1/(1-(V_R/V_{BR})^n)$ reaches $M = 2$, meaning each cycle doubles the stored charge.

    \item \textbf{$r_3 = 1.0$ (Breakdown):} Axiomatic from Axiom~4. The saturation factor reaches zero, compliance vanishes, and the topology is destroyed. In different domains this manifests as: pair production (EM), event horizon formation (gravity), quark deconfinement (nuclear), or superconducting transition (BCS). In semiconductor terms: $V_R = V_{BR}$, $M \to \infty$, the device is destroyed.
\end{itemize}

\subsection{Semiconductor Device Analogy}

The four regimes map exactly to the standard operating regions of a semiconductor p-n junction:

\begin{center}
\renewcommand{\arraystretch}{1.3}
\begin{tabularx}{\textwidth}{@{}l X X@{}}
\toprule
\textbf{Regime} & \textbf{Semiconductor} & \textbf{AVE Universal} \\
\midrule
\textbf{I: Small-Signal} &
  DC bias, linearised $g_m$, $r_\pi$, $C_\pi$ &
  $S \approx 1$, standard Maxwell/Newton \\
\textbf{II: Large-Signal} &
  Switching transient, saturation onset &
  $S(r) = \sqrt{1-r^2}$ required in full \\
\textbf{III: Avalanche} &
  Miller $M \geq 2$, carrier multiplication &
  $Q \geq 2$, energy trapping dominates \\
\textbf{IV: Breakdown} &
  $V_R = V_{BR}$, $M \to \infty$, junction failure &
  $S = 0$, topology destroyed \\
\bottomrule
\end{tabularx}
\end{center}

\noindent This is not an analogy---it is the \textbf{same operator} ($S = \sqrt{1 - (V/V_{BR})^2}$ in the semiconductor case) applied at different scales. The nuclear periodic table already classifies each element by its $V_R/V_{BR}$ ratio and Miller exponent $n$ (the proton's cinquefoil crossing number $c = 5$).

\subsection{Perturbative Expansion (Regime~I)}

In Regime~I, the saturation factor admits a convergent Taylor expansion:
\begin{equation}
    S(r) = 1 - \frac{r^2}{2} - \frac{r^4}{8} - \frac{r^6}{16} - \cdots
    \label{eq:S_taylor}
\end{equation}
To leading order, $\Delta S \approx r^2/2$. At the derived boundary $r_1 = \sqrt{2\alpha}$, $\Delta S = \alpha \approx 1/137$. This validates the use of unmodified Maxwell/Newton equations in all terrestrial laboratory experiments except those explicitly designed to approach $V_{yield}$ (e.g., asymmetrical high-voltage capacitors).

\section{Domain Control Parameter Catalog}
\label{sec:domain_catalog}

Each domain has a unique physical interpretation of the amplitude $A$ and the critical threshold $A_c$. In every case, $A_c$ is \textbf{derived from the four axioms}---it is never a fitted or empirical parameter.

\subsection{Electromagnetic (Dielectric)}

\begin{center}
\begin{tabular}{ll}
\toprule
Amplitude $A$ & Applied voltage $V$ \\
Critical $A_c$ & $V_{yield} = \sqrt{\alpha} \cdot V_{snap} = \sqrt{\alpha} \cdot m_e c^2/e \approx 43.65$ kV \\
Control parameter & $r = V / V_{yield}$ \\
\bottomrule
\end{tabular}
\end{center}

\noindent \textbf{Regime locations:}
\begin{itemize}
    \item Lab capacitor (1~kV): $r = 0.023$ --- Regime~I
    \item HV Capacitor @ 30~kV: $r = 0.687$ --- Regime~II
    \item HV Capacitor @ 43~kV: $r = 0.985$ --- Regime~III
    \item Schwinger pair production: $r = 1.0$ --- Regime~IV boundary
\end{itemize}

\subsection{Electromagnetic (Field Strength)}

\begin{center}
\begin{tabular}{ll}
\toprule
Amplitude $A$ & Local electric field $E$ \\
Critical $A_c$ & $E_{yield} = V_{yield}/\ell_{node} \approx 1.13 \times 10^{17}$ V/m \\
Control parameter & $r = E / E_{yield}$ \\
\bottomrule
\end{tabular}
\end{center}

\noindent At the lattice scale, even extreme laboratory fields ($10^6$~V/m) correspond to $r \sim 10^{-11}$, deep in Regime~I. The Schwinger critical field ($E_S = m_e^2 c^3/(e\hbar) \approx 1.32 \times 10^{18}$~V/m) corresponds to $r = E_S/E_{yield} = V_{snap}/V_{yield} = 1/\sqrt{\alpha} \approx 11.7$, deep in Regime~IV.

\subsection{Gravitational}

\begin{center}
\begin{tabular}{ll}
\toprule
Amplitude $A$ & Principal radial strain $\varepsilon_{11} = 7GM/(c^2 r)$ \\
Critical $A_c$ & Unitary strain = 1 \\
Control parameter & $r = \varepsilon_{11}$ \\
\bottomrule
\end{tabular}
\end{center}

\noindent The factor of 7 arises from the 7 compliance modes of the K4/SRS lattice ($\nu_{vac} = 2/7$). \textbf{Regime locations:}
\begin{itemize}
    \item Solar surface: $\varepsilon_{11} = 2.1 \times 10^{-6}$ --- Regime~I
    \item White dwarf: $\varepsilon_{11} \approx 3 \times 10^{-4}$ --- Regime~I
    \item Neutron star (1.4~$M_\odot$, $R = 10$~km): $\varepsilon_{11} \approx 1.46$ --- Regime~IV
    \item Black hole at $r_s = 2GM/c^2$: $\varepsilon_{11} = 7/2 = 3.5$ --- Regime~IV
\end{itemize}

\noindent \textit{The neutron star result demands attention.} A 1.4~$M_\odot$ star at $R = 10$~km has $\varepsilon_{11} = 1.46 > 1$, placing it in Regime~IV (ruptured topology). This is the AVE analog of the Buchdahl limit: no static configuration can have $r < 9GM/4c^2$ (where $\varepsilon_{11} > 63/36 = 1.75$). In AVE, the saturation boundary defines an absolute limit on gravitational compactness --- the lattice topology cannot support strain beyond unity.

\begin{examplebox}[Evaluating the Gravitational Regime of a Black Hole]
\textbf{Problem:} Calculate the dimensionless control parameter $r$ for the gravitational strain at the Event Horizon ($r_s = 2GM/c^2$) of a black hole, and determine its operating regime.

\textbf{Solution:} In the gravitational domain, the amplitude $A$ is the principal radial strain $\varepsilon_{11} = 7GM/(c^2 R_{radial})$, where the factor of 7 arises from the 7 compliance modes of the K4 lattice. The critical capacity of the lattice is unitary strain ($A_c = 1$).
At the Schwarzschild radius $R_{radial} = r_s = 2GM/c^2$:
\[ r = \frac{A}{A_c} = \frac{7GM/(c^2 \cdot (2GM/c^2))}{1} = \frac{7}{2} = 3.5 \]
Because $r = 3.5 \ge 1.0$, the local spatial metric has structurally failed. The LC network at the event horizon is in \textbf{Regime IV (Ruptured)}, meaning topological coherence is destroyed and local scalar $c_{eff} = 0$.
\end{examplebox}

\subsection{BCS / Superconducting}

\begin{center}
\begin{tabular}{ll}
\toprule
Amplitude $A$ & Temperature $T$ (or magnetic field $B$) \\
Critical $A_c$ & $T_c$ (or $B_{c0}$) \\
Control parameter & $r = T/T_c$ \\
\bottomrule
\end{tabular}
\end{center}

\noindent The BCS relation $B_c(T) = B_{c0}\sqrt{1 - (T/T_c)^2}$ IS the AVE saturation operator, with $A = T$ and $A_c = T_c$. Below $T_c$: superconducting ($S > 0$, Meissner effect active). At $T_c$: normal state ($S = 0$, topology destroyed).

\subsection{Magnetic}

\begin{center}
\begin{tabular}{ll}
\toprule
Amplitude $A$ & Applied magnetic field $B$ \\
Critical $A_c$ & $B_{snap} = m_e^2 c^2/(e\hbar) \approx 1.89 \times 10^9$ T \\
Control parameter & $r = B / B_{snap}$ \\
\bottomrule
\end{tabular}
\end{center}

\noindent Lab magnets ($\sim$10~T) have $r \sim 10^{-8}$ (Regime~I). A magnetar surface field ($\sim 10^{10}$~T) has $r \sim 5$ (Regime~IV). At $B_{snap}$, the vacuum permeability $\mu_{eff}$ reaches zero --- the vacuum becomes a perfect diamagnet (the Meissner effect of the vacuum itself).

\subsection{Nuclear}

\begin{center}
\begin{tabular}{ll}
\toprule
Amplitude $A$ & $d_{sat}/r_{sep}$ (confinement ratio) \\
Critical $A_c$ & 1 (Pauli wall) \\
Control parameter & $r = d_{sat}/r_{sep}$ \\
\bottomrule
\end{tabular}
\end{center}

\noindent When $r_{sep} > d_{sat}$: $r < 1$, the interaction potential $U(r) = -(K/r)(1 - 2\Gamma^2)$ is attractive (Regime~I--III). When $r_{sep} \leq d_{sat}$: $r \geq 1$, $\Gamma \to 1$, the potential becomes repulsive --- this is Pauli exclusion from first principles.

\subsection{Gravitational Waves}

\begin{center}
\begin{tabular}{ll}
\toprule
Amplitude $A$ & GW strain $h$ \\
Critical $A_c$ & $h_{yield} = \sqrt{\alpha} \approx 0.0854$ \\
Control parameter & $r = h / \sqrt{\alpha}$ \\
\bottomrule
\end{tabular}
\end{center}

\noindent LIGO detections ($h \sim 10^{-21}$) correspond to $r \sim 10^{-20}$ --- the most deeply linear measurement in physics. A neutron star merger at the surface ($h \sim 0.01$) reaches $r = 0.117 < \sqrt{2\alpha} = 0.121$ --- still in Regime~I. The correction $\Delta S = 0.007 < \alpha$ is sub-$\alpha$ and unresolvable. A strong merger at $h \sim 0.02$ ($r = 0.23$) enters Regime~II.

\subsection{Galactic Rotation}

\begin{center}
\begin{tabular}{ll}
\toprule
Amplitude $A$ & Newtonian gravitational acceleration $g_N$ \\
Critical $A_c$ & $a_0 \approx 1.2 \times 10^{-10}$ m/s$^2$ \\
Control parameter & $r = g_N / a_0$ \\
\bottomrule
\end{tabular}
\end{center}

\noindent \textit{Note:} The galactic domain uses the \textbf{same} universal operator as every other domain: $S(r \to 1)$ means medium compliance $\to 0$. What differs is the \textit{observational consequence}: in EM, $\varepsilon \to 0$ produces pair production (dramatic). In gravity, $G_{shear} \to 0$ produces an event horizon (dramatic). In the galactic domain, $\eta \to 0$ means the lattice drag \textit{vanishes}, leaving simple Newtonian gravity (unspectacular).

The ``dark matter problem'' \textbf{IS} the Regime~III$\to$IV phase transition. At $r = g_N/a_0 \approx 1$, the lattice mutual inductance undergoes the same phase transition as every other domain. ``Dark matter'' is the drag that exists on the Regime~I--III side of this boundary. Inner-galaxy stars at $g_N \gg a_0$ ($r \gg 1$) are in Regime~IV: the lattice is fully saturated, drag is zero, orbits are Keplerian. Outer-galaxy stars at $g_N \ll a_0$ ($r \ll 1$) are in Regime~I: the lattice is unsaturated, full drag produces $g_{eff} \to \sqrt{g_N \cdot a_0}$ (flat rotation curves).

\section{Regime-Specific Equation Sets}

For each regime, the constitutive parameters simplify:

\begin{center}
\begin{tabular}{|l|c|c|c|c|}
\hline
\textbf{Quantity} & \textbf{I (Linear)} & \textbf{II (Nonlinear)} & \textbf{III (Yield)} & \textbf{IV (Ruptured)} \\ \hline
$\varepsilon_{eff}$ & $\varepsilon_0$ & $\varepsilon_0 \sqrt{1-r^2}$ & $\to 0$ & 0 \\
$\mu_{eff}$         & $\mu_0$         & $\mu_0 \sqrt{1-r^2}$         & $\to 0$ & 0 \\
$c_{eff}$           & $c_0$           & $c_0(1-r^2)^{1/4}$          & $\to 0$ & 0 \\
$Z_{eff}$           & $Z_0$           & $Z_0 / (1-r^2)^{1/4}$       & $\to \infty$ & $\infty$ \\
$C_{eff}$           & $C_0$           & $C_0 / \sqrt{1-r^2}$        & $\to \infty$ & $\infty$ \\
$Q$                 & $\sim 1$        & $1/\sqrt{1-r^2}$             & $\gg 1$      & $\infty$ \\ \hline
\end{tabular}
\end{center}

\noindent The symmetric saturation case (both $\varepsilon$ and $\mu$ collapse equally) preserves $Z = \sqrt{\mu/\varepsilon} = Z_0$ (impedance invariant), while the asymmetric case (one collapses before the other) drives $Z$ to either 0 or $\infty$. The symmetric case governs gravity; the asymmetric case governs particle confinement (Meissner-like $\mu \to 0$ first $\Rightarrow Z \to 0$, $\Gamma \to -1$, standing wave $=$ rest mass).

\section{Cross-Domain Dimensional Analysis}
\label{sec:dimensional_analysis}

In every domain, when expressed in terms of the dimensionless control parameter $r$, the equations of motion take the same form:
\begin{resultbox}{Universal Dimensionless Master Equation}
\begin{equation}
    \frac{\partial^2 \phi}{\partial t^2} = c_0^2 (1-r^2)^{1/2} \nabla^2 \phi
\end{equation}
\end{resultbox}
where $\phi$ is the generalized displacement field (voltage, strain, temperature, GW perturbation) and $r = A/A_c$ is the domain-specific control parameter. This is a single equation governing all of physics; the domain merely specifies the physical meaning of $\phi$, $A$, and $A_c$.

\section{The Experimental Design Space}

Figure~\ref{fig:regime_design_space} maps every AVE experiment and astrophysical benchmark onto the regime diagram.

\begin{center}
\begin{tabular}{|l|c|c|c|}
\hline
\textbf{Object / Experiment} & \textbf{Domain} & \textbf{$r$} & \textbf{Regime} \\ \hline
\multicolumn{4}{|c|}{\textit{Regime I (Linear, $r < 0.1$)}} \\ \hline
LIGO GW150914              & GW       & $10^{-20}$ & I \\
Solar surface              & Gravity  & $2 \times 10^{-6}$  & I \\
Lab capacitor (1~kV)       & EM       & $0.023$    & I \\
MRI scanner (3~T)          & Magnetic & $10^{-9}$  & I \\
\hline
\multicolumn{4}{|c|}{\textit{Regime II (Nonlinear, $0.1 \leq r < 0.9$)}} \\ \hline
HV Capacitor @ 30~kV          & EM       & $0.687$    & II \\
Nb superconductor @ 4~K    & BCS      & $0.435$    & II \\
NS merger surface ($h \sim 0.01$) & GW & $0.117$    & I (sub-$\alpha$) \\
\hline
\multicolumn{4}{|c|}{\textit{Regime III (Yield, $0.9 \leq r < 1.0$)}} \\ \hline
HV Capacitor @ 43~kV          & EM       & $0.985$    & III \\
Nb @ 9.1~K ($T_c = 9.2$~K) & BCS    & $0.989$    & III \\
\hline
\multicolumn{4}{|c|}{\textit{Regime IV (Ruptured, $r \geq 1.0$)}} \\ \hline
Neutron star ($1.4 M_\odot$, 10~km) & Gravity & $1.46$ & IV \\
Black hole at $r_s$        & Gravity  & $3.50$     & IV \\
Magnetar ($10^{10}$~T)     & Magnetic & $5.29$     & IV \\
Schwinger field             & EM      & $11.7$     & IV \\
\hline
\end{tabular}
\end{center}

\begin{figure}[H]
    \centering
    \includegraphics[width=1.0\textwidth]{regime_map.png}
    \caption{\textbf{The Universal Regime Map.} The saturation factor $S(r)$ (left axis) and energy trapping quality factor $Q = 1/S$ (right axis) as functions of the dimensionless control parameter $r = A/A_c$. The four regimes are shaded: blue (Linear), yellow (Nonlinear), orange (Yield), red (Ruptured). Physical operating points from all eight domains are marked. Every experiment and astrophysical object in the AVE framework is locatable on this single diagram.}
    \label{fig:regime_design_space}
\end{figure}

\section*{Chapter Summary}
\begin{itemize}
    \item The AVE framework is formally gated by the universal dimensionless control parameter $r = A/A_c$. All physics reduce to the evaluation of the saturation operator $S(r) = \sqrt{1 - r^2}$.
    \item The regime map dictates when linear (Maxwell/Newton) approximations are valid ($r \ll 1$) and when non-linear tensor mechanics or phase transitions dominate.
    \item Seemingly disconnected extreme-physics thresholds---like Schwinger pair production, the Buchdahl compactness limit, and superconducting transition temperatures---are revealed to be the exact same macroscopic $r \to 1$ topological failure point of the vacuum hardware.
\end{itemize}

\begin{summarybox}
\begin{itemize}
    \item All physics partitions into four universal regimes classified by the dimensionless control parameter $r = A/A_c$: Linear ($r \ll 1$), Nonlinear ($r \lesssim 1$), Yield ($r \to 1$), and Ruptured ($r \geq 1$).
    \item Identical saturation mathematics produce semiconductor avalanche breakdown, black hole event horizons, and Schwinger pair-production at their respective critical thresholds.
    \item No domain analysis proceeds without first identifying its regime---the regime classification is the prerequisite gate for all AVE calculations.
\end{itemize}
\end{summarybox}

\begin{exercisebox}
\begin{enumerate}
    \item Using the first-order Taylor expansion of $S(r)$ (Eq. \ref{eq:S_taylor}), calculate the percentage magnitude error introduced by using the linear low-energy Maxwell equations to model a $30\,\text{kV}$ asymmetric vacuum capacitor where $r \approx 0.687$.
    \item According to the AVE framework, what physical property of the $\mathcal{M}_A$ lattice is destroyed when $r \ge 1.0$, and why does this conceptually eliminate the need for mathematical singularities at the center of black holes? 
\end{enumerate}
\end{exercisebox}
